% subsection content: 階差等比複合型 \, $\displaystyle a_{n+1}=pa_n+f(n)$

\subsubsection{一般型 \, $\displaystyle a_{n+1}=pa_n+f(n)$}
では次に,階差等比複合型について説明する.
この漸化式は,等比型と階差型の組み合わせからなっている.
先程の特性方程式型と同じように,等比型に帰着させることを試みよう.
\[
  a_{n+1}+g(n+1)=p\{a_n+g(n)\}
\]
となってくれれば,漸化式は等比型に帰着にする.
このような$g(n)$が存在するかどうかは分からないが,一度存在を仮定して計算を進めよう.
ここで,\,$g(n)$の次数は$f(n)$の次数に依存する.
  したがって,\,$f(n)$が$m$次式の時,\,$g(n)$は$m$次式である.
この時$g(n)$は
\[
  g(n) = x_0 + x_1 n^1 + x_2 n^2  + \cdots +x_{m} n^m =\sum_{i=0}^{m} x_i n^{i}
\]
のように書くことができる.
ここで$x_i$は適当な実数が並べられた数の集合である.
$g(n)$が求まったら,\,$b_n=a_n+g(n)$と置くと,
\begin{align*}
  b_{n+1}=pb_n
\end{align*}
のように等比数列になる.
あとは,先程の特性方程式型と同じようにして一般項を求めることができる.
\subsubsection{指数型  \, $\displaystyle a_{n+1}=pa_n+q^n$}
前節では$f(n)$が任意の関数である場合を扱った.
ここでは$f(n)$が特に指数関数である場合を扱う.
指数関数である場合にはもう一つ別の解法が存在する.
\begin{align*}
  a_{n+1}=pa_n+q^n
\end{align*}
両辺を$q^n$で割って,
\begin{align*}
   & \frac{a_{n+1}}{q^n}=\frac{p}{q}\cdot\frac{a_n}{q^{n-1}}+1 \\
   & \text{ここで,}\,b_n=\frac{a_n}{q^{n-1}}\text{と置くと,}          \\
   & b_{n+1}=\frac{p}{q}b_n+1
\end{align*}
あとは特性方程式型の漸化式として解くことになる.
